% Discussion --- plan \S3.3, "\S6".
% Sources: 03_fixed_removal.tex:266 (the author's doubt about subductive
% proliferation, carried over in impersonal register).  The rest is new.

\section{Discussion}
\label{path:sec:discussion}

\subsection{Limitations}
\label{path:sec:limitations}

Of the three candidate advantages set out in \cref{path:sec:advantages},
this chapter settles one. \emph{Limited killing} is what
\cref{path:sec:intracellular,path:sec:removal} model, and the results are
the certainty threshold, the diagonal law and
$\mathcal{R}(\vartheta)=V_\infty a^{-\vartheta}$. The other two were named
there and left, and this is where the reasons are given.

\emph{Subductive proliferation} --- a replication rate that declines during
residence in a host cell --- has its machinery already: it is a
time-inhomogeneous birth process, a special case of \cref{m:eq:meanodeinhom}
with $\lambda(t)=\beta\ee^{-\alpha t}+\gamma$, and the first-jump
probability, the mean and the deterministic trajectory are all elementary.
That material is not carried in this chapter, because it does not by itself
distinguish the two release strategies. Indeed it is not obviously an
argument for bursting at all: a rate that decays with residence time
penalises precisely the strategy that waits before releasing, so the sign of
the effect is unclear before the calculation is done, and the calculation is
not done here. The apparatus is retained as a candidate worked example for
\cref{m:sec:inhomogeneous}, where the general result it specialises already
lives.

\emph{Detection avoidance} is named and not pursued because this chapter
has no model of detection. It would need a hazard that responds to the rate
of release rather than to its size, which is the opposite dependence from
the one modelled here, and it is the mechanism most often invoked for HIV,
where release is observed to be clustered and where the responding immune
compartment has models of its own \cite{conway2018modeling}.

Four further restrictions are worth stating.

The replicator--suppressor process of \cref{path:def:rs} lets the pathogen
grow without bound once the store is spent. Real intracellular capacity is
finite, and a full model would terminate the run at rupture rather than let
$X_t$ diverge. Nothing in \cref{path:prop:certainty} or
\cref{path:prop:diagonal} depends on the late-time behaviour --- both are
statements about which absorbing set is reached --- but the survival maps of
\cref{path:sec:rsnumerics} should be read as probabilities of avoiding
clearance, not as predictions of intracellular load.

The granule store is never replenished (\cref{path:rem:noreplenish}). This
is the assumption that makes depletion permanent, and it is what turns a
larger cohort into a lasting advantage rather than a temporary one. With
replenishment at any positive rate the certainty threshold would soften into
a race between refill and consumption.

The removal budget is fixed, and \cref{path:fig:logremoval,path:fig:capped}
show what is being idealised: a slowly saturating casualty count is replaced by a
step. The replacement is what makes \cref{path:eq:R0} available in closed
form, and the honest claim is the one made in \cref{path:sec:playingfield},
that the qualitative dynamics are expected to be shared with the logarithmic
case rather than that the two agree; which of a model's predictions survive
such a substitution is not a question that answers itself
\cite{deboer2012our}. It also fixes $\vartheta$ as an integer
(\cref{path:rem:integer}), which is right for a count of killed particles
and wrong for an average over deliveries.

Finally, \cref{path:prop:reversal} compares single batches. It matches two
laws at one delivery each, exactly as \cref{p:thm:flood} does, and
\cref{path:rem:scope} sets out what that does and does not license. In
particular it is not a statement about a whole infection, and the
population-level consequences of a threshold clearance are not derived
anywhere in this thesis.

\subsection{Open problems}
\label{path:sec:open}

\begin{enumerate}[leftmargin=1.6em,itemsep=0.25em]
\item \textbf{The chained process at $\mu>0$.} The joint laws of the rupture
times and sizes $(t_k,r_k)$ are known in closed form only for $\mu=0$.
\Cref{dist:sec:chained-scope} states why the decomposition fails once a
lineage can vanish internally, and states the problem first; this chapter
inherits it, because \cref{path:sec:chainedruptures} is the model it
concerns.
\item \textbf{Off the diagonal.} \Cref{path:prop:diagonal} gives
$\pi_{j,j}$ in closed form because level~0 of the first-step system closes
on itself. No such product form is available for $\pi_{i,j}$ with $i<j$: the
generating function in $i$ couples level $j$ to level $j-1$ with $\pi_{1,j}$
entering as an unknown fixed by a degree condition. Off-diagonal values are
computed and tabulated (\cref{path:fig:repsupp20,path:fig:repsupp10}) and
left at that.
\item \textbf{The reversal at intermediate $\varphi$.}
\Cref{p:sec:open-problems} asks where on the partial-release continuum of
\cref{p:sec:partial} the flooding advantage of \cref{p:thm:flood} survives.
\Cref{path:prop:reversal} answers the same question from the other side, at
the two endpoints, and finds the criterion reversed. What happens at
intermediate $\varphi$ under a threshold is not known, and the two
sub-questions are plainly the same problem viewed from opposite ends.
\item \textbf{Clustered deliveries and a recovering field.} If the killer
field recovers on a timescale comparable with the interval between
deliveries, then \cref{path:rem:reset} fails and the toll paid at successive
deliveries is correlated. \Cref{path:sec:spectrum} notes that
self-excitation is the one mechanism in \cref{p:sec:reach} for which this
matters, and the question of whether reduced tolls can compensate for extra
deliveries needs a model that carries the recovery timescale, which this one
does not.
\end{enumerate}

\subsection{Connections forward}
\label{path:sec:forward}

The selection gradient identified in \cref{path:sec:spectrum} is a
statement about variation between release phenotypes, and \ChEvo{} is where
variability itself becomes the object of study rather than a property of a
fixed mechanism. The \yp{} account of \cref{path:sec:ypestis} is the piece
this chapter contributes to \fwd{conclusion}{the concluding chapter}, whose
subject is the evolved strategy of that organism: what is offered here is
one testable number, the threshold budget $\vartheta^{\ast}$, and the claim
that the switch between the intracellular and extracellular phases should
occur where the effector density crosses it.

\FloatBarrier
